\chapter*{Abstract}
\addcontentsline{toc}{chapter}{Abstract}
Most formal theories of information begin after information has been constituted: Shannon's theory starts from a probability space of messages, semantic theories from propositions, information algebra from a set of information pieces. This monograph studies the step before: how a raw presentation---a measurement, a record, a model activation---becomes information \emph{about} a target, \emph{for} a question.

The central object is the quotient of a \emph{contrast domain} of alternatives induced by a family of views. In the deterministic kernel (Part~I) this quotient is a partition; questions are partitions of the same domain, and a question is answerable exactly when it factors through the quotient. View sets and quotients are linked by a Galois connection, and the induced quotient has a universal property. The later parts relax the kernel's standing assumptions one at a time: admissible, non-canonical registration (Part~II); noise, graded answerability and the coupling of marginal descriptions (Part~III); anchoring, mis-anchoring and the fidelity of attribution (Part~IV); descent and obstruction to gluing (Part~V); composite readings and their interaction (Parts~VI--VII); corruption, closure and control (Part~VIII); the enriched quantitative evaluation (Part~IX); dynamical scenes (Part~X); and multiple domains with a global information algebra (Part~XI).

Results are finite unless stated otherwise; each theorem carries its hypotheses, and the book records, in its closing sections and in Appendix~A, which broader claims are \emph{not} established. Appendix~D collects supplementary finite theorems, including a certified-partial-answer criterion, an exact characterization of idempotent replication, and exact message passing on coordinate join trees.
\cleardoublepage
